\documentclass[12pt, letterpaper]{article}
\usepackage{spreamble}

\title{COP3530 - Quiz 2 (Unique Book Titles)}
\author{Ian Zou}

\begin{document}

\maketitle

\part*{Part I - Algorithm Design}

This is a naive solution that is implemented without C++ Sets and Tables.
\bigskip

\begin{minted}[
        linenos,
        frame=lines,
        framesep=2mm,
        baselinestretch=1.2,
        breaklines,
        mathescape,
]
{cpp}
#include <vector>
#include <string>
#include <iostream>

bool is_in_vec(const std::vector<std::string>& vec, const std::string& query) {
    for (size_t i = 0; i < vec.size(); ++i) {
        if (vec.at(i) == query) {
            return true;
        }
    }

    return false;
}
std::vector<std::string> unique_strings(const std::vector<std::vector<std::string>>& strings) {
    std::vector<std::string> uniques = {};
    for (size_t row = 0; row < strings.size(); ++row) {
        for (size_t col = 0; col < strings[row].size(); ++col) {
            if (is_in_vec(uniques, strings[row][col])) {
                continue;
            }

            uniques.push_back(strings[row][col]);
        }
    }

    return uniques;
}
\end{minted}

\part*{Part II - Algorithm Analysis}

The function \texttt{unique\_strings} iterates through the entire 2-D array, checking whether each string is in \texttt{uniques} with the function \texttt{is\_in\_vec}, which will iterate through the provided vector to check whether the string is within that vector. If \texttt{is\_in\_vec} returns false, the unique string will be pushed to the back of the \texttt{uniques} vector.
\medskip

In the worst case, we assume that all strings have $ c $ characters and all strings are unique. We must iterate through the characters of each string in the 2-D array, and for each string, check whether there is a string in the \texttt{uniques} array that is the same as the current string in the 2-D array. Since each string is unique, the worst-case time complexity of \texttt{is\_in\_vec} is $ \Omega (n * c) $, where $ n $ is the number of strings in the array. There are $ n * m $ strings in the 2-D array, where $ n $ is the number of rows and $ m $ is the number of columns. Thus, the worst-case time complexity of \texttt{unique\_strings}, therefore, is $ \Omega (n * m * c) $.
\medskip

In the worst case, since we assume all strings have $ c $ characters annd all are strings are unique, the \texttt{uniques} vector will be of length $ n * m $, where $ n $ is the number of rows and $ m $ is the number of columns. The worst-case auxiliary space complexity of \texttt{unique\_strings}, therefore, is $ \Omega (n * m * c) $.

\end{document}
